\subsection{Von freien zu wechselwirkenden Quantenfeldern}
\label{subsec:freie_wechselwirkende_quantenfelder}

In Kapitel~\ref{chap:relativistisches_elektron} wurde der Übergang von der
relativistischen Beschreibung eines einzelnen Elektrons zur
Quantenfeldtheorie\index{Quantenfeldtheorie} vollzogen.
Damit ist geklärt, in welchem modernen theoretischen Rahmen Elektronen,
Positronen und Photonen beschrieben werden.

Eine entscheidende Frage ist jedoch noch offen.
Die bisher betrachteten Felder erklären zunächst, welche Zustände und
Teilchenarten auftreten können.
Sie legen aber noch nicht fest, wie das Elektronenfeld und das
elektromagnetische Feld miteinander gekoppelt sind.

Die zentrale Leitfrage dieses Abschnitts lautet daher:

\begin{quote}
	\textbf{Wie wird aus einer Theorie freier Quantenfelder eine Theorie
	ihrer Wechselwirkung?}
\end{quote}

Diese Frage führt zur Quantenelektrodynamik\index{Quantenelektrodynamik},
kurz QED\@.
Sie beschreibt nicht nur die freien Felder, sondern vor allem die
elektromagnetische Wechselwirkung zwischen ihnen.

\subsubsection*{Freie Felder}

Ein freies Quantenfeld\index{Quantenfeld!freies} wird zunächst ohne
Wechselwirkung mit anderen Feldern betrachtet.
Seine zeitliche Entwicklung folgt allein seiner eigenen
Bewegungsgleichung.

Für das freie Dirac-Feld bedeutet dies:
Die Dynamik wird durch Masse, Impuls und Spinstruktur des Elektrons
bestimmt, aber noch nicht durch ein äußeres elektromagnetisches Feld.

Auch das elektromagnetische Feld kann zunächst frei betrachtet werden.
Dann breitet es sich ohne elektrische Ladungen und Ströme aus, die es
erzeugen, verändern oder streuen.

Der Begriff \emph{frei} bezeichnet somit keine neue Art von Feld.
Er kennzeichnet eine idealisierte Situation, in der der
Wechselwirkungsterm\index{Wechselwirkungsterm} der Theorie weggelassen
wird.

\begin{DidacticBox}[Was bedeutet ein freies Feld?]
	\label{box:freies_feld_ohne_wechselwirkung}

	Ein freies Feld ist weder unbewegt noch physikalisch wirkungslos.
	Es kann Energie und Impuls tragen und sich durch den Raum ausbreiten.

	\emph{Frei} bedeutet lediglich, dass seine Wechselwirkung mit anderen
	Feldern in der betrachteten Gleichung nicht berücksichtigt wird.

	Freie Felder bilden den mathematischen Ausgangspunkt. Reale
	Streu-, Emissions- und Absorptionsprozesse erfordern zusätzlich eine
	Kopplung zwischen den Feldern.
\end{DidacticBox}

Die Beschreibung durch freie Felder ist dennoch von großer praktischer
Bedeutung.
Vor einer Wechselwirkung können einfallende Teilchen häufig
näherungsweise als freie Zustände behandelt werden.
Nach der Wechselwirkung entfernen sich die erzeugten oder gestreuten
Teilchen wieder voneinander und können erneut näherungsweise durch freie
Zustände beschrieben werden.

Die eigentliche Wechselwirkung findet zwischen diesen beiden
Grenzsituationen statt.
Gerade diese Trennung bildet später die Grundlage für die Berechnung von
Streuprozessen.

\subsubsection*{Die freie Dirac-Gleichung in kovarianter Form}

In Kapitel VII wurde die Dirac-Gleichung in einer Form behandelt, die
der zeitabhängigen Schrödingergleichung ähnelt.
Für den Übergang zur Quantenelektrodynamik ist eine kompaktere
Schreibweise zweckmäßig.

Die freie Dirac-Gleichung\index{Dirac-Gleichung!kovariante Form} kann
kovariant geschrieben werden als

\begin{equation}
	\left(
	\mathrm{i}\hbar c\,\gamma^\mu\partial_\mu
	-
	m_{\mathrm e}c^2
	\right)
	\psi(x)
	=
	0.
	\label{eq:kovariante_diracgleichung_qed}
\end{equation}

Die Umformung der Dirac-Gleichung von der Hamiltonform in die kovariante
Schreibweise wird in
Anhang~\ref{app:diracgleichung_kovariante_form} schrittweise hergeleitet.

Das Symbol \(\psi(x)\) bezeichnet hier zunächst eine klassische, noch
nicht quantisierte Dirac-Feldfunktion am Raumzeitpunkt \(x\).
Diese Schreibweise dient dazu, die kovariante Struktur der
Bewegungsgleichung einzuführen.

Bei der Quantisierung wird aus dieser Feldfunktion der Feldoperator
\(\hat{\psi}(x)\), wie in Kapitel VII erläutert.
Die Ersetzung \(\psi(x)\longrightarrow\hat{\psi}(x)\) verändert jedoch
nicht die hier eingeführte kovariante Form der Dirac-Gleichung.
Sie erweitert die Bedeutung des Feldes, nicht seine relativistische
mathematische Struktur.

Die kompakte Gleichung enthält dieselbe freie relativistische Dynamik wie
die bereits bekannte Hamiltonform der Dirac-Gleichung.
Sie ordnet Raum und Zeit jedoch in einer gemeinsamen
vierdimensionalen Schreibweise an.

Dazu fasst man Zeit und Raum in einer
Viererkoordinate\index{Viererkoordinate} zusammen:

\begin{equation}
	x^\mu
	=
	\left(ct,\mathbf r\right),
	\qquad
	\mu=0,1,2,3.
	\label{eq:viererkoordinate_qed}
\end{equation}

Der Index \(\mu=0\) bezeichnet die zeitliche Komponente.
Die Werte \(\mu=1,2,3\) bezeichnen die drei räumlichen Komponenten.

Entsprechend wird die Ableitung nach Raum und Zeit durch den
Viererableitungsoperator\index{Viererableitung}

\begin{equation}
	\partial_\mu
	=
	\frac{\partial}{\partial x^\mu}
	\label{eq:viererableitung_qed}
\end{equation}

zusammengefasst.
Die Stellung des Viererindex ist dabei nicht beliebig.
Obere und untere Indizes werden mithilfe der
Minkowski-Metrik\index{Minkowski-Metrik} ineinander überführt.
Je nach verwendeter Metrikkonvention können sich beim Anheben oder
Absenken eines Index die Vorzeichen der räumlichen Komponenten ändern.
Größen wie \(x^\mu\) und \(x_\mu\), \(A^\mu\) und \(A_\mu\) oder
\(\partial^\mu\) und \(\partial_\mu\) gehören daher jeweils zusammen,
sind aber nicht in jeder Komponente einfach identisch.
Für den weiteren Gedankengang genügt diese Unterscheidung; eine
ausführliche Tensorrechnung ist hier nicht erforderlich.

Der Ausdruck \(\partial_\mu\psi\) beschreibt damit, wie sich das
Dirac-Feld in der Raumzeit verändert.

In Gleichung~\eqref{eq:kovariante_diracgleichung_qed} tritt der Index
\(\mu\) zweimal auf.
Nach der Einsteinschen
Summenkonvention\index{Einsteinsche Summenkonvention} wird dann über
alle vier Werte
summiert:

\begin{equation}
	\gamma^\mu\partial_\mu
	=
	\gamma^0\partial_0
	+
	\gamma^1\partial_1
	+
	\gamma^2\partial_2
	+
	\gamma^3\partial_3.
	\label{eq:gamma_partial_summe_qed}
\end{equation}

Die Größen \(\gamma^\mu\) sind die bereits aus der Dirac-Theorie
stammenden Gamma-Matrizen\index{Gamma-Matrizen}.
Sie sorgen dafür, dass die vier Komponenten des Dirac-Feldes in der
richtigen Weise miteinander gekoppelt werden.

Ihre genaue Matrizenform muss hier nicht erneut hergeleitet werden.
Für den Übergang zur QED ist vor allem ihre physikalische Aufgabe
wichtig:
Sie verbinden die Spinstruktur des Elektrons mit seiner Bewegung durch
die Raumzeit.

Die Bezeichnung \emph{kovariant}\index{Lorentz-Kovarianz} bedeutet,
dass die Gleichung bei einem Wechsel zwischen gleichförmig bewegten
Bezugssystemen ihre mathematische Form behält.
Raum und Zeit erscheinen nicht als voneinander getrennte Bestandteile,
sondern als Komponenten einer gemeinsamen Raumzeitbeschreibung.

\begin{PhysicsBox}[Bedeutung der kovarianten Dirac-Gleichung]
	\label{box:bedeutung_kovariante_diracgleichung_qed}

	Die kovariante Dirac-Gleichung beschreibt die freie relativistische
	Dynamik des Elektronenfeldes:

	\begin{equation*}
		\left(
		\mathrm{i}\hbar c\,\gamma^\mu\partial_\mu
		-
		m_{\mathrm e}c^2
		\right)
		\psi
		=
		0.
	\end{equation*}

	Die Viererableitung verbindet die zeitliche und räumliche Änderung des
	Feldes. Die Gamma-Matrizen enthalten die relativistische
	Spinstruktur.

	Eine elektromagnetische Wechselwirkung ist in dieser freien Gleichung
	noch nicht enthalten.
\end{PhysicsBox}

Der letzte Punkt ist entscheidend.
Gleichung~\eqref{eq:kovariante_diracgleichung_qed} beschreibt, wie sich
das Dirac-Feld ohne elektromagnetische Kopplung entwickelt.
In ihr tritt weder ein elektrisches noch ein magnetisches Potential auf.

Für eine Theorie der Wechselwirkung muss die freie Gleichung daher
erweitert werden.

\subsubsection*{Das elektromagnetische Viererpotential}

In der klassischen Elektrodynamik können das elektrische und das
magnetische Feld mithilfe zweier Potentiale beschrieben werden:
dem skalaren elektrischen Potential \(\phi\) und dem magnetischen
Vektorpotential \(\mathbf A\).

In der relativistischen Schreibweise werden beide zum
Viererpotential\index{Viererpotential}

\begin{equation}
	A^\mu
	=
	\left(
	\frac{\phi}{c},
	\mathbf A
	\right)
	\label{eq:viererpotential_qed}
\end{equation}

zusammengefasst.

Aus den Potentialen ergeben sich die bekannten Felder:

\begin{equation}
	\mathbf E
	=
	-\nabla\phi
	-
	\frac{\partial\mathbf A}{\partial t},
	\qquad
	\mathbf B
	=
	\nabla\times\mathbf A.
	\label{eq:felder_aus_potentialen_qed}
\end{equation}

Das Viererpotential ist deshalb die natürliche relativistische Größe,
mit der das Dirac-Feld an die elektromagnetische Wechselwirkung
gekoppelt werden kann.

An dieser Stelle genügt zunächst diese strukturelle Aussage.
Die genaue Kopplung wird nicht durch ein beliebiges Hinzufügen des
Potentials festgelegt.
Sie muss mit der Relativitätstheorie, der Ladungserhaltung und der
Quantenmechanik vereinbar sein.

Zudem ist das Viererpotential nicht eindeutig bestimmt.
Verschiedene Potentiale können zu denselben messbaren elektrischen und
magnetischen Feldern führen.
Diese Freiheit wird als Eichfreiheit\index{Eichfreiheit} bezeichnet.
Sie ist kein mathematischer Mangel, sondern der Schlüssel zur Form der
elektromagnetischen Wechselwirkung.

\subsubsection*{Was eine Wechselwirkungstheorie ergänzen muss}

Bei freien Feldern können die Gleichungen für das Dirac-Feld und das
elektromagnetische Feld zunächst getrennt betrachtet werden.
In der Quantenelektrodynamik ist diese Trennung nicht mehr vollständig
möglich.

Das elektromagnetische Feld beeinflusst die zeitliche und räumliche
Entwicklung des Dirac-Feldes.
Umgekehrt ist die elektrische Ladung des Dirac-Feldes eine Quelle des
elektromagnetischen Feldes.

Beide Felder müssen daher durch einen gemeinsamen Wechselwirkungsterm
miteinander verbunden werden.

\begin{PhysicsBox}[Von den freien Feldern zur QED]
	\label{box:von_freien_feldern_zur_qed}

	Die Struktur der Quantenelektrodynamik lässt sich zunächst symbolisch
	als Summe dreier Bestandteile schreiben:

	\begin{equation*}
		\begin{aligned}
			\text{QED}
			&=
			\text{freies Dirac-Feld}
			\\
			&\quad+
			\text{freies elektromagnetisches Feld}
			\\
			&\quad+
			\text{Wechselwirkung}.
		\end{aligned}
	\end{equation*}

	Die freien Gleichungen legen fest, wie sich beide Felder ohne Kopplung
	entwickeln. Der Wechselwirkungsterm bestimmt, wie Elektronen,
	Positronen und Photonen an elektromagnetischen Prozessen beteiligt
	sind.
\end{PhysicsBox}

Damit ist die Aufgabe für den nächsten Schritt klar.
Gesucht ist keine beliebige Kopplung, sondern eine mathematisch
eindeutige Erweiterung der freien Dirac-Gleichung.

Die entscheidende Forderung lautet, dass die physikalischen Aussagen
unverändert bleiben, wenn die Phase des Dirac-Feldes an jedem
Raumzeitpunkt in geeigneter Weise verändert wird.
Diese Forderung führt zur lokalen
Eichsymmetrie\index{Eichsymmetrie!lokale}.

Im nächsten Abschnitt wird gezeigt, warum gerade diese Symmetrie das
elektromagnetische Viererpotential notwendig macht und wie daraus die
Kopplung zwischen dem Dirac-Feld und dem elektromagnetischen Feld
entsteht.
